\section{Introduction}\label{sec:intro}

A viability kernel answers the question whether a constrained system can be kept inside its constraints forever \cite{Aubin1991,Aubin2011,Blanchini1999}. Under full state observation the answer is local: the state is viable at first order when some admissible action points the drift inside the constraint set, and Nagumo's condition makes this exact for a large class of sets \cite{Nagumo1942}. Under partial observation the controller does not know the state; it knows an \emph{information atom}, the set $B$ of states compatible with what has been observed, and it must choose one action for all of them. Viability can then fail although every state in $B$ is individually controllable: the safe-action sets of the members of $B$ are non-empty but their intersection is empty. We call this the \emph{common-action obstruction}. It is the mechanism behind the delayed-observation thresholds, the shared-authority ladders and the aggregate-observation failures studied in the companion papers \cite{Abaee2026calc,Abaee2026cert,Abaee2026prob}, and it is the first-order shadow of the output-feedback viability problem.

The question of this paper is how to \emph{certify} such an obstruction: how to produce an object that proves, portably and checkably, that no common action exists. A direct proof enumerates the action set; a certificate should be small, verifiable by a single inequality, and informative about where the obstruction comes from. For control-affine plants with a compact convex action set the answer is a probability measure on the active constraints of the atom. Averaged under the measure, the constraints collapse to a single affine row, and the obstruction is exactly the statement that this row is negative throughout the action set. The equivalence is Sion's minimax theorem in the zero-sum game between the controller, who picks the common action, and an adversary, who picks the state, the constraint and the disturbance. The dual object is a measure because the adversary mixes.

\paragraph{Contributions.} The paper develops this idea into a complete static theory, a taxonomy, a benchmark law and a dynamic layer.
\begin{enumerate}[leftmargin=1.6em,itemsep=1pt]
\item \emph{Static theory} (Section~\ref{sec:static}). Theorem~\ref{thm:dual} states the duality with its exactness: an information atom is obstructed if and only if a measure certificate exists, and the certificate margin equals the negative of the lower value of the game. Remark~\ref{rem:isaacs} identifies the measure dual with the lower value and separates it from the upper value, which is the full-information quantity. Theorem~\ref{thm:helly} bounds the support of a certificate by $\min(\dim\aff U,\rank)+1$ atoms through Helly's theorem---not $k+2$ as Carath\'eodory alone would give---and exhibits, for every $k$, a plant on which the bound is attained. Proposition~\ref{prop:nonconvex} settles the non-convex case: soundness is unconditional, and a certificate exists exactly when the convexified atom is obstructed, so the gap between $U$ and $\conv U$ is the only source of incompleteness. Theorem~\ref{thm:polyhedral} recovers the Farkas form for polyhedral $U$ and shows that the facet multipliers locate the obstruction on specific constraints of the action set.
\item \emph{Taxonomy} (Section~\ref{sec:taxonomy}). Three tests---physical, epistemic, institutional---are defined on the safe-action sets of an atom (Proposition~\ref{prop:taxonomy}). They are tests, not a partition: the institutional test is relative to a reference action set and overlaps with either of the other two. The certificate's support and its facet multipliers read the tests off.
\item \emph{Benchmark law} (Section~\ref{sec:benchmark}). On a two-stock harvesting model with an aggregate observation and a demand floor $D$, Proposition~\ref{prop:benchmark} gives the fibre margin $\varepsilon^\ast(Y)=(27-5Y)/50$, the unique certificate $(\tfr12,\tfr12)$ and its Farkas pair, and Proposition~\ref{prop:floor} the piecewise critical law $Y^\ast(D)=\min(\tfr{77}5-5D,\tfr{69}5)$ whose two branches are, respectively, an epistemic--institutional and a physical obstruction. Each unit of demand floor costs five units of aggregate stock.
\item \emph{Dynamic layer} (Section~\ref{sec:dynamic}). Finite review trees with action-dependent observation cells are defined (Definition~\ref{def:tree}); Theorem~\ref{thm:envelope} proves that the node recursion is the supremum over adapted policy trees, with the tower and martingale identities along any fixed tree, and Corollary~\ref{cor:prefix} that committed prefixes inherit node values. Example~\ref{ex:posterior} shows that posterior-only evaluation is not an envelope. Proposition~\ref{prop:robust} states the certified weak duality of the robust value over a class of couplings and its mixed-policy closure; Example~\ref{ex:sweep} tabulates envelope, open-loop and both robust values on a two-window instance, exhibits an \emph{interior} worst coupling and a dynamically certified obstruction. Proposition~\ref{prop:integrated} is the integrated certificate for finite-horizon trajectories with adapted controls; on the three-branch instance of the companion paper \cite{Abaee2026cert} it reproduces the review threshold $\tau^\ast=\tfr7{50}$ in one line.
\item \emph{Verification} (Section~\ref{sec:verification}). Every number, witness and count is regenerated from the definitions by an exact rational script; the soundness direction, the parity gap and the benchmark algebra are checked in Lean~4. A scope box separates what is proved from what is only computed on instances.
\end{enumerate}

\paragraph{What is new relative to the viability and robust-optimisation literature.} Output-feedback viability has sufficient conditions through set-valued observers and guaranteed viability kernels \cite{Aubin2011,Cardaliaguet2007}; the present dual gives a necessary-and-sufficient first-order criterion with a certificate. The measure certificate is the semi-infinite programming dual \cite{HettichKortanek1993,AndersonNash1987} specialised to affine rows; its Helly-type sparsity is the deterministic counterpart of the support-constraint count of the scenario approach \cite{CalafioreCampi2006,CampiGaratti2008}, where $k+1$ also appears. The dynamic envelope is the finite-horizon dynamic programme of a partially observed problem \cite{Astrom1965,BertsekasShreve1978} written on review trees, which is what makes the tree-supremum identity, the committed-prefix corollary and the posterior-only counterexample elementary; the robust value over couplings is a static-information design problem in the sense of Bergemann and Morris \cite{BergemannMorris2019} with the worst coupling as the designer. The adjustability question of two-stage robust optimisation \cite{WeiZhang2024}, when adaptivity adds value, is the open-loop gap of Section~\ref{sec:dynamic}.

\paragraph{Reading guide.} Section~\ref{sec:setting} fixes the plant, the information atoms and the row abstraction. Sections~\ref{sec:static}--\ref{sec:dynamic} contain the results; proofs that are short are given in place, the longer ones in the supplementary material, which also contains the complete instance records, the random study, the Lean scope list and a symbol table with sign conventions. Section~\ref{sec:verification} describes the verification record, Section~\ref{sec:related} places the results and lists open problems.
